\chapter{Troubleshooting and Data Safety}

\section{Startup problems}

\begin{description}
  \item[Splash screen stuck on ``Extracting the CAD engine\ldots{}''] The first start, or the first start after an update,
        needs to unpack about 1 GB of runtime. Keep the screen awake and make sure at least 2 GB of free space remains.
        If there is still no progress after 3 minutes, force-stop the app and open it again; if it keeps failing,
        uninstall and reinstall (documents on the storage card are not affected).
  \item[Red text ``CAD engine failed to start: \ldots{}''] Note the error message after the red text. Common causes:
        not enough storage, a device ABI mismatch (this APK supports arm64 only), or a broken system WebView/runtime.
        Try rebooting the device first; if it keeps failing, report the error message.
  \item[The app closes immediately after starting] Check whether a repackaged build whose signature conflicts with
        the official one is installed; uninstall it and reinstall the package from the appropriate channel.
\end{description}

\section{Feature failures}

Red text or an exclamation mark in the feature tree means the feature cannot be computed. Work through these steps in order:

\begin{enumerate}
  \item Double-click the feature to open its dialog and read the error at the top (for example ``Cannot find the geometry on the datum feature'' or ``The geometry does not exist on the base feature'');
  \item Check its \textbf{references}: is the profile closed, does it self-intersect, do the referenced faces and edges still exist (changing an earlier feature may have invalidated the references);
  \item Use \ui{Suppress} to skip it temporarily and check whether later features depend on it;
  \item Go back to the sketch and work through the \ui{Constraints} list (conflicting, redundant, malformed), or delete the extra line segments;
  \item Once fixed, long-press the feature and choose \ui{Recompute object}.
\end{enumerate}

\section{Sketch problems}

\begin{description}
  \item[The status bar keeps showing ``Under-constrained: N degrees of freedom''] You are still N constraints short.
        Common omissions: a point not fixed near the origin, a missing horizontal or vertical constraint, a missing dimension.
        Go through the \ui{Constraints} list one by one.
  \item[Conflicting] Two constraints contradict each other (for example you gave both horizontal and vertical distances
        and an angle as well). Delete one of them. Conflicting constraints are shown in red, and the \ui{Constraints}
        panel can show only the conflicting or redundant ones.
  \item[Redundant] A constraint is repeated (the same edge constrained twice); delete the extra one.
  \item[The sketch flips] When a sketch is under-constrained, dragging geometry or changing a dimension can flip the whole
        shape. Constrain the sketch to ``fully constrained'' and add a lock to the key points.
  \item[Driving or reference dimensions] A dimension can be switched to ``reference (driven)'': it only shows the measured
        value and does not take part in constraining, which is useful for checking.
\end{description}

\section{Files and export}

\begin{description}
  \item[A file opens blank or reports an error] Make sure the format is supported; for \file{.FCStd}, save with a similar
        FreeCAD version. A complex STEP/IGES downloaded from the web can take a long time to import.
  \item[No solids after exporting STEP] The model may contain only meshes or surfaces; check first that it is a valid solid with ``Part → Check geometry''.
  \item[The STL surface is very rough] Reduce the deviation when exporting (0.01--0.05 mm), or refine the mesh first.
  \item[The exported PDF is a blank page] The drawing page has no views yet, or the view scale is so large that it runs off the sheet; check the page and view properties.
\end{description}

\section{Crashes and recovery}

MobileCAD has crash recovery: it saves a recovery copy every 60 seconds. After a crash, reopen the app and it offers to restore. Recommendations:
\begin{itemize}
  \item save manually at important milestones with \menupath{File → Save};
  \item use \menupath{File → Save a copy} to take version snapshots;
  \item if the recovery prompt does not appear, check \file{files/recovery/} in the app's private directory (root or a debug build required).
\end{itemize}

\section{Update problems}

\begin{description}
  \item[``Update failed: checksum mismatch''] The download is corrupt; just retry. The app verifies SHA-256.
  \item[``Update not installed''] You must allow MobileCAD to install unknown apps: follow the prompt into system settings, grant permission, and come back.
  \item[The F-Droid build does not offer updates] That is normal: the F-Droid build updates through F-Droid, and the app shows ``updates come from the installation source''.
\end{description}

\section{Performance and memory}

\begin{itemize}
  \item In large assemblies, hide unrelated objects first; on complex models, avoid keeping several live-preview dialogs open at once;
  \item An FEM solve saturates the CPU, so watch for heat and throttling; start with a few thousand elements;
  \item when free storage drops below 2 GB the system may reclaim background apps, so save before switching apps.
\end{itemize}

\section{Feedback and help}

\begin{itemize}
  \item Project repository: \url{https://github.com/LucasCoraca/mobilecad} (in an issue, include the device model, Android version,
        MobileCAD version, FreeCAD commit, reproduction steps, and a screenshot or screen recording);
  \item Problems in FreeCAD itself (geometry kernel, algorithms) can also be reported to the upstream FreeCAD project;
  \item Corrections to this manual: open an issue in the repository and tag it ``manual''.
\end{itemize}
